% Siden den industrielle revolution har fossile brændstoffer været menneskets primære energikilde.
% Forbrænding af fossile brændstoffer frigiver store mængder drivhusgasser, som bidrager til global opvarmning og klimaforandringer. 
% Derfor er der stigende fokus på alternative energikilder som sol, vind og kernekraft.

% Kernekraft udnyttes i dag primært gennem fission, hvor tunge atomkerner sammen spaltes og frigiver energi. 
% Fission har det negative biprodukt, radioaktivt affald, som medfører en række sikkerheds/miljømæssige udfordringer.

% Fusion, hvor lette atomkerner sammensmeltes, betragtes som en lovende fremtidig energikilde. 
% Processen kan potentielt levere store mængder energi hvorved affaldet har meget højere halveringstid.

% For at opnå fusion opvarmes en kilde af atomer til ekstremt høj temperatur, hvor stoffet overgår 
% til plasma, en ioniseret gas bestående af frie elektroner og ioner. 
% I moderne fusionsanlæg som tokamaker fastholdes plasmaet ved hjælp af stærke elektro magnetiskefelter,
% så det ikke kommer i kontakt med tokamakkens vægge, hvilket medfører energitab og ustabilitet.

En central udfordring inden for fusionsforskning er at forstå og beskrive plasmaets dynamik, 
særligt i randzonen, scarpe-off layer (SOL), hvor turbulens og transportfænomener har stor betydning for energitab og stabilitet.
Disse processer er modelleret med koblede partielle differentialligninger (PDE'er). 
Projektet fokuserer på hot edge SOL electrostatic (HESEL) modellen.\cite{PDE_til_plasma}

Til at simulere HESEL-modellen bruges numerisk finite difference. Simulering er et centralt matematisk værktøj i moderne plasmafysik.
En generel ulempe ved simuleringsarbejde er, at det er beregningsmæssigt tungt og bestemt af begyndelsesbetingelserne.
Plasmaets randdynamik er stærkt ikke-lineær og følsom over for små ændringer i begyndelsesbetingelserne, 
hvilket gør det vanskeligt at forudsige disse processer nøjagtigt uden at køre omfattende simuleringer.\cite{SciML_til_plasma}

I dette projekt undersøges det, hvordan Scientific Machine Learning (SciML) kan anvendes til at modellere og simulere 
randdynamik i plasmaet med fokus på HESEL-modellen.
SciML har vist sig at kunne generalisere godt til simuleringsdata samt at overholde de fysiske love. 
Ved at træne en SciML-model på en lille mængde simuleringsdata vil man kunne bygge en statistisk model, 
som fortæller, hvordan systemet vil opføre sig.
En veltrænet SciML-model vil kunne generalisere HESEL-modellen og fortælle, hvordan systemet vil opføre sig 
givet nogle interpolerede begyndelsesbetingelser, punkter og tidsskridt uden at skulle køre omfattende simuleringer.

Projektet vil først fokusere på ablation af HESEL-modellen 
til træning af state of the art (SOTA) SciML-metoder, specifikt Physics Informed Neural Networks (PINN).
Herefter implementeres og trænes disse modeller på det fulde HESEL-ligningssystem og med forskellige begyndelsesbetingelser.

Målet er at undersøge, om SciML kan bidrage til mere fleksibel og effektiv modellering af 
plasmaets randdynamik.